\begin{abstract}
Large language models (LLMs) can reduce the effort of drafting behavior-driven development (BDD) scenarios, but valid syntax does not guarantee that business intent is preserved. We evaluate GPT-5.4-mini on Connextra-style user stories reverse-engineered from 261 expert-written feature files in Sylius, Apache Fineract and Diaspora, using one same-project example per prompt. We measure parser compliance with the behave Gherkin parser and similarity to the reference with skeleton cosine similarity (all-MiniLM-L6-v2). Excluding the three prompt exemplars ($N = 258$), 255 files (98.8\%, 95\% CI [96.6\%, 99.8\%]) are valid Gherkin, well above the 80\% threshold ($p < 0.001$). The median similarity is 0.853, but a one-sided Wilcoxon test does not show that it exceeds the 0.85 threshold ($p = 0.474$). Only the e-commerce project exceeds it; the finance and social-network projects lie significantly below it, and the conclusion changes at a threshold of 0.84. Among 60 high-similarity files, 25 contain a business-logic defect in AI-drafted, human-reviewed labels (agreement with an independent rater: $\kappa = 0.57$), and seeded defects barely change the score. The similarity result also depends on the prompt: without an example the median falls to 0.819, and with a different exemplar per project it rises to 0.875, significantly above 0.85, while parser compliance stays at 98.8--99.6\%. GPT-5.4-mini is therefore a drafting aid whose business logic needs human review, not an autonomous acceptance-test generator, and embedding similarity alone cannot certify generated Gherkin.

\keywords{Large language models \and Gherkin \and Behavior-driven development \and Acceptance testing \and Software testing \and GPT-5.4-mini}
\end{abstract}
